\documentclass[12pt, a4paper]{ctexart}

% ==================== 宏包 ====================
\usepackage[top=2.5cm, bottom=2.5cm, left=2.8cm, right=2.8cm]{geometry}
\usepackage{amsmath, amssymb}
\usepackage{graphicx}
\usepackage{booktabs}
\usepackage{multirow}
\usepackage{float}
\usepackage{enumitem}
\usepackage{hyperref}
\usepackage{xcolor}
\usepackage{caption}
\usepackage{subcaption}
\usepackage{fancyhdr}
\usepackage{setspace}
\usepackage{titlesec}
\usepackage{tabularx}
\usepackage{longtable}
\usepackage{array}
\usepackage{tikz}
\usepackage{pgfplots}
\pgfplotsset{compat=1.18}

% ==================== 格式设置 ====================
\hypersetup{colorlinks=true, linkcolor=blue, citecolor=blue, urlcolor=blue}
\setstretch{1.45}

\pagestyle{fancy}
\fancyhf{}
\fancyhead[C]{\small 信息与知识获取 --- 作业一调研报告}
\fancyfoot[C]{\thepage}
\renewcommand{\headrulewidth}{0.4pt}

% ==================== 标题页 ====================
\title{
    \vspace{1cm}
    {\LARGE \textbf{信息与知识获取}} \\[0.6cm]
    {\Large 作业一：智能手机传感器的国内外发展情况调研} \\[0.4cm]
    {\large 调研报告}
    \vspace{0.5cm}
}
\author{
    姓名：XXX \quad 学号：XXXXXXXXXX \\[0.2cm]
    % 成员1姓名：XXX \quad 学号：XXXXXXXXXX \\[0.2cm]
    % 成员2姓名：XXX \quad 学号：XXXXXXXXXX \\[0.2cm]
    \textbf{分工说明：}独立完成全部调研、分析与报告撰写工作
}
\date{2026 年 3 月}

\begin{document}
\maketitle
\thispagestyle{empty}
\newpage

% ==================== 摘要与关键词 ====================
\setcounter{page}{1}
\begin{abstract}
智能手机传感器是现代移动终端实现多维度环境感知的核心硬件基础。本报告围绕``智能手机传感器的国内外发展情况''这一主题，首先对智能手机中图像与光学、声音、运动与姿态、触摸与力、环境感知以及生物特征与健康六大类传感器进行了系统梳理，然后以CMOS图像传感器（CIS）作为主要调研方向，从与人类视觉感知的对比、基于内光电效应的工作原理、前照式$\to$背照式$\to$堆叠式三代架构演进等维度展开了深入分析。在产业发展方面，报告对比了国际市场（索尼约50\%份额、三星约20\%、豪威科技约12\%）与国内市场（豪威集团、思特威、格科微``三强鼎立''）的竞争格局，并分析了计算摄影与AI融合驱动的软硬协同趋势。此外，报告从健康监测、环境保护、电子废弃物治理等角度讨论了传感器技术对环境和社会可持续发展的影响，最后对多模态传感融合、AI原生传感器设计、新型交互范式和可持续性驱动创新四个方向提出了展望。
\end{abstract}

\noindent\textbf{关键词：}智能手机传感器；CMOS图像传感器；背照式CMOS；堆叠式CMOS；计算摄影；多模态传感融合；可持续发展

\newpage
\tableofcontents
\newpage

% ==================================================================
%                          正    文
% ==================================================================

\section{引言}

智能手机已经从单纯的通信工具演变为人类感知世界的延伸。这一转变的核心驱动力，正是嵌入手机内部的各类传感器。截至2026年，一部主流智能手机平均集成超过14种物理和虚拟传感器\cite{alibaba2026}，从最初的加速度计和光线感应，发展到如今涵盖图像、声音、触摸、运动、环境感知、生物识别等多维度的传感体系，使得智能手机成为了一个高度集成的便携式感知平台。

从产业角度看，全球CMOS图像传感器市场在2025年已达到105.99亿美元，预计到2034年将增长至247.38亿美元\cite{statsmarket2026}。中国传感器市场规模已超过4000亿元人民币\cite{eetChina2025}，以豪威集团、思特威、格科微为代表的国产企业正在快速崛起，在中高端市场与索尼、三星展开激烈角逐。与此同时，传感器技术正经历从单一感知向多模态融合、从硬件竞赛向软硬协同、从消费电子向汽车和医疗等领域渗透的深刻变革。

本报告首先全面梳理智能手机中各类传感器的分类与基本功能，然后以\textbf{图像传感器（CMOS Image Sensor, CIS）}作为主要调研方向，深入分析其与人类视觉感知的关系、基本工作原理，以及国内外研究与应用的发展现状，并对传感器技术对环境和社会可持续发展的影响进行讨论，最后对未来发展方向提出创新性思考与预测。图~\ref{fig:market_growth} 展示了全球CMOS图像传感器市场的增长趋势。

\begin{figure}[H]
\centering
\begin{tikzpicture}
\begin{axis}[
    ybar,
    bar width=16pt,
    xlabel={年份},
    ylabel={市场规模（亿美元）},
    symbolic x coords={2020, 2022, 2024, 2025, 2028, 2031, 2034},
    xtick=data,
    nodes near coords,
    every node near coord/.append style={font=\footnotesize},
    nodes near coords align={vertical},
    ymin=0,
    ymax=280,
    width=0.82\textwidth,
    height=7cm,
    enlarge x limits=0.1,
]
\addplot[fill=blue!40] coordinates {(2020, 68) (2022, 82) (2024, 98) (2025, 106) (2028, 140) (2031, 186) (2034, 247)};
\end{axis}
\end{tikzpicture}
\caption{全球CMOS图像传感器市场规模及预测（2020--2034年）\cite{statsmarket2026}}
\label{fig:market_growth}
\end{figure}


\section{智能手机传感器全面概述}

\subsection{传感器分类与基本功能}

智能手机中的传感器按照感知对象的不同，可以划分为以下六大类：

\textbf{（1）图像与光学感知类传感器。}包括CMOS图像传感器（主摄像头、超广角、长焦、前置摄像头）、飞行时间（Time-of-Flight, ToF）深度传感器、结构光传感器（如iPhone的Face ID）和激光雷达（LiDAR）扫描仪。这类传感器赋予手机"看"的能力，是手机摄影、人脸识别、增强现实（AR）等功能的硬件基础。

\textbf{（2）声音感知类传感器。}包括MEMS麦克风阵列和骨传导传感器。现代智能手机通常配备2--3个MEMS麦克风，通过波束成形（beamforming）技术实现定向收音和环境噪声抑制，赋予手机"听"的能力。

\textbf{（3）运动与姿态感知类传感器。}包括三轴加速度计、三轴陀螺仪、磁力计（电子罗盘）和气压计。加速度计和陀螺仪通常集成为六轴惯性测量单元（IMU），为步数计算、屏幕旋转、导航定位和运动游戏提供核心数据。磁力计检测地球磁场方向以实现指南针功能，气压计通过测量大气压力推算海拔高度。

\textbf{（4）触摸与力感知类传感器。}包括电容式触摸屏传感器、压力传感器（如Apple的3D Touch/Haptic Touch）和超声波指纹传感器。电容式触摸屏通过检测人体手指接近时引起的电容变化来定位触摸位置，支持多点触控；超声波指纹传感器利用超声波穿透OLED屏幕采集指纹3D信息，实现屏下指纹解锁。

\textbf{（5）环境感知类传感器。}包括环境光传感器（ALS）、接近传感器、温度传感器和湿度传感器。环境光传感器测量周围光照强度以自动调节屏幕亮度，接近传感器在通话时检测手机与面部的距离以关闭屏幕。自2025年3月起，所有新发布的智能手机均标配温度传感器\cite{alibaba2026}。

\textbf{（6）生物特征与健康感知类传感器。}包括光电容积脉搏波（PPG）传感器（用于心率和血氧测量）、红外温度传感器和光谱传感器。然而，目前仅32\%的Android设备和68\%的iPhone在心率和血氧测量方面达到临床级精度标准\cite{alibaba2026}，该领域仍有较大提升空间。

\subsection{主要传感器功能汇总}

表~\ref{tab:sensor_overview} 对智能手机中的主要传感器进行了系统性汇总。

\begin{table}[H]
\centering
\caption{智能手机主要传感器汇总}
\label{tab:sensor_overview}
\small
\begin{tabular}{llll}
\toprule
\textbf{类别} & \textbf{传感器名称} & \textbf{核心功能} & \textbf{对应人类感知} \\
\midrule
\multirow{4}{*}{图像/光学} & CMOS图像传感器 & 拍照、录像 & 视觉 \\
 & ToF深度传感器 & 3D深度感知 & 立体视觉 \\
 & 结构光传感器 & 人脸3D建模 & 立体视觉 \\
 & LiDAR扫描仪 & 空间扫描与建模 & 空间感知 \\
\midrule
声音 & MEMS麦克风 & 录音、语音识别 & 听觉 \\
\midrule
\multirow{4}{*}{运动/姿态} & 加速度计 & 运动检测、计步 & 前庭觉 \\
 & 陀螺仪 & 旋转检测、防抖 & 前庭觉 \\
 & 磁力计 & 方向感知 & （无直接对应） \\
 & 气压计 & 海拔测量 & （无直接对应） \\
\midrule
\multirow{2}{*}{触摸/力} & 电容触摸屏 & 触摸交互 & 触觉 \\
 & 超声波指纹 & 身份识别 & 触觉 \\
\midrule
\multirow{3}{*}{环境} & 环境光传感器 & 亮度调节 & 视觉（明暗适应） \\
 & 接近传感器 & 距离检测 & 触觉/空间感 \\
 & 温度传感器 & 环境温度测量 & 温度觉 \\
\midrule
\multirow{2}{*}{生物/健康} & PPG传感器 & 心率、血氧检测 & （无直接对应） \\
 & 红外温度传感器 & 体温测量 & 温度觉 \\
\bottomrule
\end{tabular}
\end{table}


\section{图像传感器深度调研}

在所有智能手机传感器中，CMOS图像传感器无疑是技术含量最高、市场规模最大、发展最为迅猛的一类。它不仅是手机摄影功能的核心，更是计算机视觉、增强现实、自动驾驶等人工智能应用的基础数据来源。本节以图像传感器为主要调研方向，从其与人类视觉的关系、工作原理和技术架构演进三个维度展开深入分析。

\subsection{图像传感器与人类视觉感知的对比}

CMOS图像传感器的设计在很大程度上借鉴了人类视觉系统的工作原理，但两者之间也存在本质性差异。表~\ref{tab:human_vs_cmos} 从多个维度对比了人眼与CMOS图像传感器的感知特性。

\begin{table}[H]
\centering
\caption{人类视觉系统与CMOS图像传感器对比}
\label{tab:human_vs_cmos}
\small
\begin{tabular}{lll}
\toprule
\textbf{对比维度} & \textbf{人类视觉} & \textbf{CMOS图像传感器} \\
\midrule
感光元件 & 视网膜：1.2亿视杆细胞 + & 像素阵列：数百万至2亿 \\
 & 600万视锥细胞 & 个光电二极管 \\
\midrule
色彩感知 & 三种视锥细胞（L/M/S） & Bayer滤色器阵列 \\
 & 感知红/绿/蓝三原色 & （RGGB排列） \\
\midrule
动态范围 & $\sim$120 dB（明暗适应） & 60--110 dB（HDR模式下） \\
\midrule
分辨力 & 中央凹区约0.3角分 & 取决于像素密度 \\
 & （等效约576百万像素） & 和镜头光学质量 \\
\midrule
暗光适应 & 视杆细胞暗适应 & 大像素、高ISO、 \\
 & （需约30分钟） & 多帧降噪（即时） \\
\midrule
对焦方式 & 睫状肌调节晶状体 & 相位检测自动对焦 \\
 & 曲率（连续变焦） & （PDAF）/激光对焦 \\
\midrule
信号传输 & 视神经→视觉皮层 & 模拟信号→ADC→ \\
 & （并行处理） & ISP数字处理 \\
\midrule
光谱范围 & 380--780 nm（可见光） & 350--1000 nm \\
 & & （含近红外） \\
\bottomrule
\end{tabular}
\end{table}

\textbf{相似之处：}从宏观架构看，CMOS图像传感器与人眼存在深层的功能类比。人眼的视网膜相当于传感器的像素阵列，视锥/视杆细胞相当于光电二极管，三种视锥细胞的分工（感受红、绿、蓝三色光）与Bayer滤色器的RGGB排列异曲同工。人脑视觉皮层的图像处理功能则对应ISP（图像信号处理器）的去马赛克、白平衡、降噪等计算过程。

\textbf{关键差异：}人类视觉系统最显著的优势在于其极高的动态范围（约120 dB）和非线性响应特性。人眼能够同时感知强烈阳光下和阴暗角落中的细节，这是因为视网膜的光感受器具有对数响应特性。相比之下，传统CMOS传感器的动态范围仅有60--70 dB，即使采用最新的LOFIC（Lateral Overflow Integration Capacitor）技术，也仅能达到约110 dB\cite{adhocnews2025}。此外，人类视觉具有注意力机制——中央凹的高分辨率区域仅覆盖约2度视角，大脑通过快速眼动（saccade）在不同兴趣区域之间切换，构建出完整的高分辨率场景感知。这一机制启发了计算摄影中的"注意力引导"超分辨率重建算法。

\textbf{传感器的超越：}在某些维度上，CMOS传感器已经超越了人类视觉。其光谱响应可延伸至近红外区域（$>$780 nm），使得夜视、红外人脸识别等功能成为可能。传感器能够以每秒数十帧甚至数百帧的速度持续采集图像，而人眼的时间分辨率约为60 Hz。此外，传感器的输出可以被精确量化和存储，为计算机视觉和机器学习提供了客观、可重复的数据基础。

\subsection{CMOS图像传感器工作原理}

CMOS图像传感器的核心工作原理基于半导体的\textbf{内光电效应}：当入射光子的能量大于硅材料的禁带宽度（约1.12 eV，对应波长$\sim$1100 nm）时，光子被硅吸收并激发出电子-空穴对，从而将光信号转换为电信号。

一个完整的成像周期包括以下四个阶段：

\textbf{（1）复位阶段。}在曝光开始前，每个像素中的光电二极管通过复位晶体管被充电至参考电压，清除上一帧残留的电荷，确保每次曝光从一致的基线开始。

\textbf{（2）曝光与光电转换阶段。}复位完成后，光线通过片上微透镜汇聚，经过彩色滤光片（Bayer滤色器按RGGB排列，使每个像素仅记录红、绿、蓝中的一种颜色信息），到达光电二极管。光子被硅吸收，产生的光生电子在光电二极管中积累。积累的电子数与入射光强和曝光时间成正比：
\begin{equation}
    N_e = \eta \cdot \Phi \cdot A \cdot t_{\text{exp}}
\end{equation}
其中 $N_e$ 为光生电子数，$\eta$ 为量子效率，$\Phi$ 为入射光子通量密度，$A$ 为像素面积，$t_{\text{exp}}$ 为曝光时间。

\textbf{（3）电荷-电压转换阶段。}曝光结束后，光电二极管中积累的电荷通过传输晶体管转移至浮动扩散节点（Floating Diffusion, FD），FD的电容 $C_{FD}$ 将电荷转换为电压信号：
\begin{equation}
    V_{FD} = V_{\text{reset}} - \frac{N_e \cdot q}{C_{FD}}
\end{equation}
其中 $q$ 为电子电荷（$1.6 \times 10^{-19}$ C）。该电压由像素内的源极跟随器放大器缓冲后输出。

\textbf{（4）读出与数字化阶段。}模拟电压信号经列级模数转换器（ADC）转换为数字信号，再由ISP进行去马赛克（将单色Bayer数据插值为全彩RGB图像）、白平衡、伽马校正、降噪、锐化等处理，最终输出完整的数字图像。

\subsection{技术架构演进：从前照式到堆叠式}

CMOS图像传感器的架构经历了三代重要演进，每一代都带来了感光性能的质的飞跃（见表~\ref{tab:cmos_evolution}）。

\begin{table}[H]
\centering
\caption{CMOS图像传感器三代架构对比}
\label{tab:cmos_evolution}
\small
\begin{tabular}{llll}
\toprule
\textbf{特性} & \textbf{前照式（FSI）} & \textbf{背照式（BSI）} & \textbf{堆叠式（Stacked）} \\
\midrule
结构特点 & 金属层在光电 & 光电二极管朝向 & 像素层与电路层 \\
 & 二极管上方 & 入射光一侧 & 分层制造后堆叠 \\
\midrule
光线利用率 & $\sim$70\% & $\sim$90\%+ & $\sim$90\%+ \\
\midrule
灵敏度提升 & 基准 & +30\%$\sim$50\% & +30\%$\sim$50\% \\
\midrule
电路复杂度 & 受限 & 略有提升 & 大幅提升 \\
\midrule
制程优化 & 统一制程 & 统一制程 & 可差异化制程 \\
\midrule
读出速度 & 慢 & 中等 & 快（减少果冻效应） \\
\midrule
代表产品 & 早期手机摄像头 & Sony Exmor R & Sony Exmor RS \\
 & & （2009年量产） & （2012年量产） \\
\midrule
典型应用 & iPhone 4及更早 & iPhone 4S$\sim$6 & iPhone 6S至今 \\
\bottomrule
\end{tabular}
\end{table}

\textbf{前照式（FSI）CMOS}是最早的架构形式。光线需要先穿过金属布线层才能到达底部的光电二极管。由于金属层对可见光的反射率高达约90\%，实际到达光电二极管的光线仅为入射光的约70\%\cite{tencent_cmos}，导致低光环境下噪声明显。

\textbf{背照式（BSI）CMOS}通过将像素结构翻转，使光线直接照射到光电二极管，绕过了金属布线层的遮挡。索尼于2009年率先实现背照式CMOS的量产化（品牌名Exmor R），并于2011年应用于iPhone 4S\cite{csdn_cmos}。BSI技术使传感器灵敏度提升了30\%--50\%，暗光成像质量显著改善，目前已成为智能手机图像传感器的标准配置。

\textbf{堆叠式（Stacked）CMOS}进一步将像素区域和信号处理电路分别在不同的硅片上制造——像素层可采用优化感光性能的65 nm制程，信号处理层则可采用更先进的45 nm甚至28 nm制程——然后通过硅通孔（TSV）或混合键合（Hybrid Bonding）技术将两层堆叠在一起\cite{tencent_cmos}。这种分离制造策略带来了多重优势：处理电路的晶体管数量可以翻倍，从而实现更快的读出速度（减少果冻效应）、片上HDR合成、高帧率拍摄（如1600万像素下120 fps）等高级功能。


\section{图像传感器国内外发展现状}

\subsection{国际发展现状}

全球智能手机CMOS图像传感器市场由三大巨头主导：

\textbf{索尼半导体解决方案公司（Sony Semiconductor Solutions）}长期占据全球CIS市场约50\%的份额，是高端手机图像传感器领域的绝对领导者。索尼的IMX系列传感器广泛应用于苹果iPhone、三星Galaxy S/Note旗舰系列、华为和小米等品牌的高端机型。在技术创新方面，索尼持续引领行业：2012年推出堆叠式CMOS（Exmor RS）；2019年推出首颗手机用4800万像素传感器（IMX586）；2023年推出双层晶体管像素技术（2-Layer Transistor Pixel），将饱和信号量提升约一倍；2025年最新的IMX775采用RGB-IR双模式成像技术，在940 nm近红外波段量子效率达35\%，动态范围达110 dB\cite{statsmarket2026}。

\textbf{三星半导体（Samsung LSI）}是全球第二大CIS供应商，市场份额约20\%。三星在高像素领域走在前列，于2021年率先推出1.08亿像素传感器ISOCELL HM3，2023年进一步推出2亿像素ISOCELL HP3，并在2025年将双2亿像素配置引入旗舰产品\cite{adhocnews2025}。三星在像素合并技术（Tetrapixel/Nonapixel）、ISOCELL Plus像素隔离技术方面拥有深厚积累。

\textbf{豪威科技（OmniVision Technologies）}全球排名第三，市场份额约10\%--15\%。作为最早进入CIS领域的企业之一（1995年成立于美国硅谷），豪威科技于2019年被中国韦尔股份收购后，成为中国CIS领域的领军企业。豪威在中低端市场占据优势地位，近年来积极向高端市场进军，其OV50H系列5000万像素传感器和最新的旗舰级产品正逐步获得高端手机品牌的设计采用（Design Win）\cite{adhocnews2025}。图~\ref{fig:market_share} 直观地展示了主要厂商的市场份额分布。

\begin{figure}[H]
\centering
\begin{tikzpicture}
\begin{axis}[
    xbar,
    bar width=14pt,
    xlabel={市场份额（\%）},
    ytick={1,2,3,4},
    yticklabels={其他厂商, 豪威科技, 三星半导体, 索尼半导体},
    xmin=0,
    xmax=62,
    ymin=0.3,
    ymax=4.7,
    width=0.7\textwidth,
    height=5.5cm,
    nodes near coords,
    nodes near coords align={horizontal},
    every node near coord/.append style={font=\small\bfseries},
]
\addplot[fill=blue!40] coordinates {(18, 1) (12, 2) (20, 3) (50, 4)};
\end{axis}
\end{tikzpicture}
\caption{全球智能手机CIS市场份额分布（2025年）\cite{statsmarket2026}}
\label{fig:market_share}
\end{figure}

\subsection{国内发展现状}

中国手机CIS市场形成了\textbf{"三强鼎立"}的竞争格局\cite{qqCIS2025}：

\textbf{豪威集团（原韦尔股份）}以1602亿元市值稳居行业首位，2024年图像传感器业务收入达191.90亿元，占其主营收入的74.76\%。2025年5月，韦尔股份正式更名为"豪威集团"，以突显图像传感器核心业务定位\cite{sensorexpert2025}。豪威集团具备全栈技术布局能力，拥有TheiaCel（超级像素技术）、PureCel Plus-S（高灵敏度像素技术）和Nyxel（近红外增强技术）等多项核心技术平台，产品线覆盖从手机到汽车、安防、医疗等全场景\cite{laoyaoba2025}。

\textbf{思特威（SmartSens Technology）}表现强劲，2024年营收首次突破50亿元并实现翻倍增长，被业界誉为"黑马"。2025年上半年智能手机CIS业务收入达32.91亿元，同比增长269.05\%，其5000万像素旗舰级CIS芯片出货量激增\cite{qqCIS2025}。思特威以Lofic HDR 2.0和SmartGS-2 Plus等差异化技术平台为核心竞争力，强调全场景技术覆盖\cite{laoyaoba2025}。

\textbf{格科微（GalaxyCore）}2024年出货量达18.26亿颗，全球市场份额约23\%，按出货量计排名全球第二\cite{qqCIS2025}。格科微的竞争策略以"单芯片集成"和极致成本控制见长，主打5000万像素以下中低端产品。其GalaxyCell 2.0平台通过工艺创新和自有产能优化，在性价比方面具有显著优势\cite{laoyaoba2025}。

在技术发展方向上，三家企业形成了以下共识\cite{laoyaoba2025}：（1）\textbf{5000万像素}已成为中端产品标配；（2）\textbf{2亿像素}产品加速进入验证与量产阶段；（3）\textbf{单次曝光HDR}技术成为主流，动态范围接近110 dB。

\subsection{国内外对比分析}

表~\ref{tab:cn_vs_intl} 从多个维度对比了国内外CIS产业的发展状况。

\begin{table}[H]
\centering
\caption{国内外CIS产业对比分析}
\label{tab:cn_vs_intl}
\small
\begin{tabular}{lll}
\toprule
\textbf{对比维度} & \textbf{国际（索尼、三星）} & \textbf{国内（豪威、思特威、格科微）} \\
\midrule
技术成熟度 & 深厚积累（30+年） & 快速追赶（10$\sim$15年） \\
高端市场份额 & 占据主导地位（$>$70\%） & 加速渗透（$<$30\%） \\
像素工艺 & 领先（双层晶体管、 & 跟进中（LOFIC、 \\
 & 堆叠式混合键合） & 单次曝光HDR） \\
制程节点 & 28 nm$\sim$22 nm & 40 nm$\sim$28 nm \\
核心优势 & 技术创新引领 & 成本控制与快速响应 \\
产业链完整度 & IDM模式（自有晶圆厂） & Fabless为主 \\
新兴领域布局 & 汽车、医疗、工业 & 汽车电子、AI眼镜、 \\
 & & AR/VR \\
\bottomrule
\end{tabular}
\end{table}

总体而言，国际厂商在高端CIS技术和旗舰市场上仍保持明显领先优势，但国内企业的追赶速度令人瞩目。特别是在中端市场（5000万像素级别），国产CIS的性能已逐渐接近国际水平，部分指标甚至实现超越。在汽车电子这一战略性新兴领域，思特威的汽车CIS业务2025年上半年同比增长107.97\%\cite{qqCIS2025}，展现出强劲的增长势头。


\section{计算摄影与AI融合：传感器的软硬协同}

2025--2026年，智能手机影像领域最显著的变革趋势是\textbf{从硬件参数竞赛向计算摄影的软硬协同方向转型}\cite{itbear2026}。传感器不再仅仅是"光的捕获者"，而是与AI算法深度融合的"智能感知节点"。

\subsection{硬件策略的理性回归}

经历了"大底竞赛"（1英寸传感器）和"像素战争"（2亿像素）之后，2025--2026年的旗舰手机传感器尺寸从1英寸向1/1.3$\sim$1/1.4英寸回调，在感光面积与机身厚度、功耗之间寻求更优的平衡点\cite{itbear2026}。行业的焦点从单一主摄的极致参数转向\textbf{全焦段均衡发展}——确保超广角、主摄和长焦三颗镜头都具备足够的成像品质。

\subsection{AI算法的深度介入}

在传感器硬件趋于成熟的背景下，AI算法成为拉开产品差距的关键要素：

\textbf{（1）多帧合成与计算HDR。}传感器在极短时间内连续拍摄多帧不同曝光的图像，AI算法将这些帧对齐、融合并去噪，生成一张动态范围远超单帧极限的照片。Google的HDRNet和Apple的Deep Fusion均属此类技术。

\textbf{（2）AI超分辨率重建。}利用深度学习模型（如Real-ESRGAN）从低分辨率的传感器原始数据中重建出高分辨率细节，使中端传感器也能输出接近旗舰水平的图像质量。

\textbf{（3）语义理解与场景自适应。}AI系统能够识别拍摄场景（如人像、风景、食物、文档等）并自动调整传感器的曝光参数、白平衡和后处理风格，这类似于人脑视觉皮层对不同场景进行差异化处理的机制。

\textbf{（4）端侧AI推理。}Android 15和iOS 19均优先采用端侧（on-device）传感器融合处理而非云端计算，带宽节省高达73\%，同时降低延迟并增强隐私保护\cite{alibaba2026}。

\subsection{传感器-算法协同设计}

值得关注的新趋势是\textbf{传感器硬件与算法的协同设计}（Sensor-Algorithm Co-design）。传感器厂商不再仅提供原始图像数据，而是在芯片级别嵌入AI加速能力，例如：片上HDR合成（在传感器内部完成多曝光融合，降低ISP负担）；智能ROI（Region of Interest）检测（传感器自动识别画面中的关键区域，优先分配读出带宽）；事件驱动感知（仅在像素值发生变化时才触发读出，大幅降低功耗和数据量）。这一趋势标志着图像传感器正从被动的"光电转换器"向主动的"智能视觉前端"演进。


\section{其他重要传感器发展趋势}

\subsection{运动传感器：从消费娱乐到专业级应用}

六轴IMU（加速度计+陀螺仪）是智能手机中数据输出最为可靠的传感器之一\cite{alibaba2026}。其应用场景已从早期的屏幕旋转和游戏交互，扩展至惯性导航、运动健康分析、地震预警等专业级场景。MEMS陀螺仪的零偏稳定性从早期的$>$10$°$/h提升至当前的$<$1$°$/h，结合磁力计和气压计的多传感器融合定位精度已能达到室内亚米级。

\subsection{声学传感器：MEMS麦克风的智能化}

MEMS麦克风正经历从单纯声音采集器向"智能听觉前端"的转型。最新一代MEMS麦克风集成了本地化的AI信号处理能力，可在传感器内部完成关键词检测（如"Hey Siri"唤醒）、声源定位和环境噪声分类，无需唤醒主处理器即可实现低功耗的始终聆听（Always-On Listening）功能。

\subsection{生物传感器：向临床级医疗迈进}

光电容积脉搏波（PPG）传感器正从运动健康监测向无创血糖检测、连续血压监测等临床级需求演进。电子皮肤技术和柔性传感器的突破使得高保真生物信号采集成为可能\cite{innosensor2025}。这一趋势与联合国可持续发展目标SDG 3（良好健康与福祉）高度契合，有望为全球数十亿智能手机用户提供低成本的基础健康监测服务。

\subsection{新兴传感器技术}

多项前沿技术正在加速从实验室走向商业化\cite{innosensor2025}：（1）\textbf{量子传感器}利用量子态的极端敏感性实现超高精度测量，可将磁场检测灵敏度提升数个数量级；（2）\textbf{光子集成电路（PIC）传感器}将光学器件集成在芯片上，实现片上光谱分析，有望让智能手机具备物质成分检测能力；（3）\textbf{太赫兹传感器}随6G通信技术的发展，"通信+传感"一体化将使手机兼具精密成像和安全检测功能；（4）\textbf{QD-on-CMOS短波红外（SWIR）传感器}通过量子点技术降低SWIR传感器成本，扩展手机在农业、食品安全检测等领域的应用。


\section{传感器对环境和社会可持续发展的影响}

\subsection{正面影响}

\textbf{（1）促进全民健康监测。}智能手机传感器使基础健康监测（心率、血氧、运动量）从医疗机构延伸到每个人的口袋，有助于缩小健康服务的城乡差距和国际差距，推动实现SDG 3（良好健康与福祉）。特别是在医疗资源匮乏的发展中国家和偏远地区，手机PPG传感器可作为基础筛查工具，以极低成本实现心血管疾病的早期预警。

\textbf{（2）赋能环境监测与灾害预警。}加速度计可用于地震预警网络（如Google的Android Earthquake Alerts已覆盖全球多个地震频发国家），气压计可辅助气象预报，图像传感器配合AI可实现空气质量可视化评估、森林火灾早期检测等环保应用，助力SDG 13（气候行动）。

\textbf{（3）推动智慧农业与食品安全。}光谱传感器和近红外成像技术的手机化应用，使农民能够通过手机检测土壤养分、作物健康状况和果蔬新鲜度，降低农药和化肥的过度使用，促进SDG 2（零饥饿）和SDG 12（负责任消费和生产）。

\textbf{（4）促进信息无障碍。}图像传感器配合OCR和AI技术，使视障人士能够通过手机"阅读"实体文本；声学传感器与语音识别技术结合，帮助听障人士实现实时语音转文字，推动SDG 10（减少不平等）。

\subsection{环境挑战}

\textbf{（1）电子废弃物问题。}智能手机中包含超过60种元素材料\cite{stcn2025}，其中多种为稀有金属和重金属。一块手机电池中的镉可污染6万升水\cite{ndrc2022}。中国每年产生4亿部以上废旧手机，但资源回收利用率仅约4\%，约54.2\%被消费者闲置\cite{stcn2025}。

\textbf{（2）稀土与关键材料依赖。}传感器制造依赖多种稀土和关键矿物（如铟用于触摸屏的ITO导电膜、钽用于电容器、钕用于磁力计中的稀土磁铁），这些矿物的开采过程伴随着严重的环境破坏和碳排放。

\textbf{（3）制造过程的碳足迹。}半导体制造是高能耗、高水耗和高化学品消耗的工业过程。一颗先进制程的CMOS图像传感器芯片的制造涉及数百道光刻、蚀刻和沉积工序，其碳足迹不容忽视。

\subsection{可持续发展路径}

\textbf{（1）绿色制造。}传感器行业正加速向低碳化制造转型\cite{innosensor2025}：开发无铅压电材料替代传统含铅PZT（锆钛酸铅）材料；采用可回收封装设计降低电子废弃物处理难度；向12英寸晶圆制造过渡以提高单位面积产出，降低资源浪费。

\textbf{（2）循环经济。}废旧手机被称为"城市矿产"——每吨废旧手机可提炼约200克黄金、2200克白银和100千克铜\cite{stcn2025}。推动手机传感器的模块化设计和可拆卸结构，有助于提高关键材料的回收率。新修订的《废弃电器电子产品处理污染控制技术规范》已于2026年3月1日起实施，将智能设备纳入管控范围\cite{jingji2026}。

\textbf{（3）低功耗设计。}传感器的本地化智能处理趋势（在传感器内部完成特征提取和初级判断，而非将原始数据上传云端）可显著降低数据传输的能耗。事件驱动式传感器仅在有意义的信号变化发生时才触发工作，进一步降低功耗。

\textbf{（4）传感器赋能碳中和。}传感器本身也可以成为可持续发展的工具——用于氢能泄漏监测、碳足迹追踪、工业能效优化和可再生能源系统的运行监控\cite{innosensor2025}。


\section{未来发展方向展望与创新思考}

基于本次调研的分析，笔者对智能手机传感器的未来发展提出以下预测和创新性建议：

\subsection{多模态传感融合将成为核心范式}

未来的智能手机将不再依赖单一传感器完成特定任务，而是通过\textbf{多模态传感融合}实现更高层次的环境感知。全球传感器融合市场预计到2032年将达到319.1亿美元，年复合增长率18.65\%\cite{innosensor2025}。例如，将图像传感器的视觉信息、IMU的运动数据、麦克风的声学信号和气压计的环境数据融合在一起，手机可以构建出对周围世界的多维度理解，从"感知单一物理量"跃升至"理解复杂场景"。

\subsection{从感知到认知：传感器的AI原生化}

笔者认为，下一代智能手机传感器将实现\textbf{"AI原生"设计}——AI不再是传感器数据的后处理工具，而是从传感器架构设计之初就被深度嵌入。具体而言：（1）传感器芯片内集成神经网络加速器，实现片上推理；（2）传感器具备自适应学习能力，根据用户使用习惯动态优化工作参数；（3）传感器之间形成去中心化的协同感知网络，无需经过主处理器即可完成跨模态信息交互。这一演进路径将使手机从"工具"向"具备初级认知能力的伙伴"转变。

\subsection{传感器驱动的下一代交互方式}

触摸屏并非人机交互的终极形态。笔者预测以下新型交互范式将在3--5年内进入主流：

\textbf{（1）空间手势交互。}结合ToF深度传感器和毫米波雷达（如Google Soli芯片的后继技术），手机将能够识别用户的空中手势，实现非接触式操控——在烹饪、手术等双手不便触碰屏幕的场景下价值巨大。

\textbf{（2）脑机接口的初步探索。}随着EEG（脑电图）传感器微型化技术的进步，未来的智能手机或配套可穿戴设备可能集成简化版脑电传感器，实现基于注意力和情绪状态的辅助交互。

\textbf{（3）嗅觉与味觉传感。}电子鼻（基于金属氧化物半导体气体传感器阵列）和电子舌技术虽尚处早期阶段，但一旦实现足够的微型化和低成本化，将使智能手机真正覆盖人类的全部五感，开启气味识别、食品安全快检等全新应用场景。

\subsection{可持续性驱动的传感器创新}

在全球碳中和目标的驱动下，笔者建议传感器行业加速以下方向的研发：（1）基于生物可降解材料的一次性传感器，用于环境监测后可自然分解；（2）能量收集型传感器（从光、热、振动中获取微能量），实现自供电运行，减少对电池的依赖；（3）传感器的"数字护照"——在芯片中记录其全生命周期的材料组成和碳足迹信息，为循环回收提供数据支撑。


\section{总结}

本报告围绕"智能手机传感器的国内外发展情况"这一主题，首先全面梳理了智能手机中六大类传感器的基本功能与感知特性，然后以CMOS图像传感器为主要调研方向，从与人类视觉的对比、内光电效应工作原理、前照式$\to$背照式$\to$堆叠式的三代架构演进、以及国内外产业竞争格局四个维度展开了深入分析。

调研发现：（1）在国际层面，索尼以约50\%的市场份额和持续的技术创新（双层晶体管像素、RGB-IR双模成像等）保持领先，三星在高像素领域追赶激进；（2）在国内层面，以豪威集团（1602亿元市值）、思特威（营收翻倍增长）和格科微（全球出货量第二）为代表的"三强"正在快速缩小与国际巨头的技术差距，在5000万像素级别已具备竞争力；（3）2025--2026年的核心趋势是计算摄影与AI融合驱动的"软硬协同"，传感器正从被动的光电转换器向智能视觉前端演进；（4）传感器技术在促进健康监测、环境保护和信息无障碍方面具有巨大的正面社会价值，但电子废弃物、稀土依赖和制造碳足迹等环境挑战也亟需产业界和政策制定者共同应对。

面向未来，笔者提出了多模态传感融合、AI原生传感器设计、新型交互范式（空间手势、脑机接口、电子鼻/舌）和可持续性驱动创新（生物可降解传感器、能量收集传感器、数字护照）四个方向的创新思考。智能手机传感器技术的持续进化，将深刻改变人类感知世界、与世界交互的方式，其影响远超消费电子领域本身。


\begin{thebibliography}{20}

\bibitem{alibaba2026}
    Master Smartphone Sensor Data Collection: A 2026 Guide.
    \url{https://www.alibaba.com/product-insights/master-smartphone-sensor-data-collection-a-2026-guide.html}, 2026.

\bibitem{innosensor2025}
    最新预测：2026年传感器的十大趋势.
    \url{http://www.innosensor.com/newsDetail_725.html}, 2025.

\bibitem{statsmarket2026}
    CMOS Image Sensor for Smartphone Market: Size, Share, Volume 2026 to 2033.
    Stats Market Research, 2026.

\bibitem{adhocnews2025}
    The Battle for Smartphone Camera Dominance: Can OmniVision Challenge the Leaders?
    \url{https://www.ad-hoc-news.de/boerse/news/}, 2025.

\bibitem{itbear2026}
    手机影像新趋势：2025至2026年技术迭代与体验升级的深度变革.
    ITBear比尔科技, 2026.

\bibitem{sensorexpert2025}
    1600亿市值国产传感器巨头改名：证券简称从"韦尔股份"变更为"豪威集团".
    传感器专家网, 2025.

\bibitem{qqCIS2025}
    格科微、思特威、韦尔股份（豪威）：国产CIS三强谁在掉队？
    腾讯新闻, 2025.

\bibitem{laoyaoba2025}
    2025年A股CIS龙头三国杀：格科微、思特威、豪威集团的差异化竞速.
    集微网, 2025.

\bibitem{eetChina2025}
    中国传感器企业2025上半年财报解析：新兴场景崛起与智能化升级.
    电子工程专辑, 2025.

\bibitem{tencent_cmos}
    堆栈式 CMOS、背照式 CMOS 和传统 CMOS 传感器的区别.
    腾讯云开发者社区.
    \url{https://cloud.tencent.com.cn/developer/article/1934252}.

\bibitem{csdn_cmos}
    CMOS图像传感器架构的演变.
    CSDN博客.
    \url{https://blog.csdn.net/qq_21842097/article/details/125438789}.

\bibitem{stcn2025}
    旧手机，如何物尽其用（大数据观察）.
    证券时报网.
    \url{https://www.stcn.com/article/detail/1084003.html}.

\bibitem{ndrc2022}
    为废旧手机找好"归宿".
    国家发展和改革委员会, 2022.
    \url{https://www.ndrc.gov.cn/fggz/hjyzy/zyzhlyhxhjj/202205/t20220523_1325139.html}.

\bibitem{jingji2026}
    废电器回收新规2026年3月1日实施.
    21经济网, 2026.

\bibitem{ebiotrade2025}
    废弃智能手机在水生系统中元素浸出行为的一致性特征及其环境风险研究.
    生物通, 2025.

\bibitem{researchmarkets2025}
    CMOS Image Sensors Market -- Global Forecast 2025--2030.
    Research and Markets, 2025.

\bibitem{adenekan2024}
    Adenekan R A G, Yoshida K T, Okamura A M, et al.
    A Comparative Analysis of Smartphone and Standard Tools for Touch Perception Assessment.
    \textit{IEEE Transactions on Haptics}, 2024.

\bibitem{mdpi2024}
    Comparative Assessment of Multimodal Sensor Data Quality Collected Using Android and iOS Smartphones in Real-World Settings.
    \textit{Sensors}, 2024, 24(19): 6246.

\bibitem{baidu_cmos}
    CMOS图像传感器.
    百度百科.
    \url{https://baike.baidu.com/item/CMOS图像传感器}.

\bibitem{gansuan2025}
    2026年传感器的十大趋势.
    感算商城, 2025.
    \url{https://m.gansuan.com/tech/sharing/2877528897366462464.html}.

\end{thebibliography}

\end{document}
